

\chapter{Mở đầu}
\label{chap:mo-dau}

Chương đầu tiên của khóa luận trình bày bối cảnh và lý do chọn đề tài, cùng với các nội dung liên quan. Mục đầu tiên mô tả vấn đề đang hiện hữu của bài toán số hóa trong lĩnh vực học tập. Hai mục tiếp theo, báo cáo sẽ trình bày tổng quan giải pháp đề xuất và các đóng góp của khóa luận. Mục cuối cùng trình bày sơ bộ về cấu trúc của báo cáo và nội dung chính của từng chương.

\section{Đặt vấn đề}

Tác vụ số hóa tài liệu học là một công việc quan trọng trong lĩnh vực giáo dục, đặc biệt trong thời đại công nghệ xuất hiện trong mọi lĩnh vực của đời sống. Đặc điểm của ngành này là có lượng dữ liệu dồi dào và vô cùng đa dạng như: sách giáo khoa, đề thi và bài giảng. Bên cạnh việc chuyển đổi tài liệu giáo dục từ dạng vật lý sang dữ liệu số, việc sử dụng chúng một cách hữu ích cũng là một bài toán thiết yếu. Bằng cách tạo ra các bài kiểm tra, đặc biệt là bài kiểm tra trắc nghiệm, giáo viên vừa có thể đánh giá học sinh một cách khách quan và tự động, vừa có thể tận dụng nguồn tài nguyên số này.

Tuy nhiên, một vấn đề lớn được đặt ra là: quá trình thiết kế ra một bộ câu hỏi trắc nghiệm, vừa đảm bảo cả chất lượng và số lượng, yêu cầu lượng thời gian và tài nguyên đáng kể từ giáo viên. Chưa kể, bài kiểm tra trắc nghiệm sau đó cần được kiểm tra và rà soát để đảm bảo các tiêu chuẩn giáo dục, như thang đo Bloom. Hai vấn đề kỹ thuật chính mà giáo viên phải đối mặt khi cần chuyển đổi học liệu sang một bài kiểm tra trắc nghiệm là:

\begin{itemize}
\item \textbf{Khó khăn trong việc số hóa tài liệu từ ảnh chụp:} Các tài liệu giáo dục thực tế thường ở dạng ảnh in, trang sách, hoặc ảnh chụp từ điện thoại với chất lượng thấp (ảnh nhòe, mờ, độ phân giải thấp). Không chỉ thế, nội dung của chúng còn bao gồm các dạng dữ liệu phức tạp như bảng biểu, đồ thị và công thức toán, gây khó khăn cho các công cụ truyền thống phổ biến như Tesseract. Hậu quả là văn bản được trích xuất ra có độ chính xác kém, thiếu nội dung và phá vỡ cấu trúc của văn bản; dẫn tới ảnh hưởng nghiêm trọng đến công việc sinh bài kiểm tra.


\item \textbf{Khó khăn trong việc kiểm soát chất lượng câu hỏi:} Các hướng tiếp cận truyền thống đối với bài toán này thường sử dụng các quy tắc lập trình sẵn (rule-based) để tạo ra các câu hỏi trắc nghiệm có mức độ nhận thức đơn giản. Chủ yếu là các câu hỏi về thông tin bề mặt, không yêu cầu suy luận và phân tích sâu, rất rập khuôn và thiếu tự nhiên. Một hướng đi mới là ứng dụng khả năng hiểu văn bản của các mô hình ngôn ngữ lớn, để từ văn bản đầu vào có thể tạo ra các câu hỏi liên quan với chất lượng sư phạm cao. Tuy nhiên, chi phí để sử dụng các mô hình này là rất lớn, đòi hỏi tài nguyên tính toán cao, đặc biệt đối với lượng học liệu khổng lồ của ngành giáo dục. Không chỉ vậy, ngay cả các LLM cũng khó tránh được việc sinh ra các câu hỏi kém chất lượng, ảo giác, nếu không có một quy trình kiểm tra chặt chẽ.

\end{itemize}

Từ những thách thức trên, một vấn đề lớn được đặt ra là: Làm thế nào có thể phát triển một hệ thống tự động số hóa tài liệu, tạo ra các bài kiểm tra trắc nghiệm có chất lượng cao với chi phí thấp, có thể triển khai trên hạ tầng có sẵn và bảo mật?

\section{Giải pháp đề xuất}

Nhằm giải quyết bài toán trên, khóa luận đề xuất giải pháp kết hợp các mô hình ngôn ngữ cỡ nhỏ mã nguồn mở (sLLM - Small Large Language Mode) với kiến trúc đa tác tử (MAS). Các mô hình được sử dụng có số lượng tham số dưới 15 tỷ, điều này giúp chúng có thể chạy trên hạ tầng sẵn có với chi phí thấp. Bên cạnh đó, kiến trúc tác tử thiết lập một quy trình tạo bài kiểm tra chặt chẽ từ bước thiết kế câu hỏi đến bước đánh giá và rà soát câu hỏi vừa sinh. Quy trình đề xuất gồm hai giai đoạn chính:

 \item \textbf{Giai đoạn trích xuất văn bản bằng VLM nhỏ gọn:} Nhận thấy được những hạn chế cố hữu của các công cụ truyền thống như testeract, giải pháp sử dụng các mô hình VLM chuyên dụng cỡ nhỏ như DeepSeek OCR v2, PaddleOCR-VL. Các mô hình này đã chứng minh được khả năng trích xuất văn bản một cách chính xác từ ảnh đầu vào, giúp đảm bảo chất lượng của giai đoạn tạo bài kiểm tra.

   \item \textbf{Giai đoạn sinh câu hỏi trắc nghiệm bằng kiến trúc đa tác tử:} Kiến trúc này tách bài toán thành nhiều bài toán nhỏ hơn được xử lý bởi một tác tử chuyên biệt. Hệ thống Lang Graph được sử dụng để điều phối các tác tử:
    \begin{itemize}
\item \textbf{Tác tử Sinh câu hỏi:} Tạo mới hoặc tạo lại các câu hỏi trắc nghiệm với đầu vào là tài liệu và mức độ nhận thức hệ thống yêu cầu.
        \item \textbf{Tác tử Phân loại:} Phân loại và kiểm tra các câu hỏi được sinh ra có đảm bảo đúng mức độ nhận thức theo thang đo Bloom hay không.
                \item \textbf{Tác tử Học sinh:} Tác tử này đóng vai là học sinh và sẽ giải các câu hỏi vừa được sinh ra mà không xem đáp án. Quá trình này giúp phát hiện ra các câu hỏi bị đa nghĩa, bẫy logic bất hợp lý, hoặc phương án nhiễu quá dễ nhận biết.

        \item \textbf{Tác tử Hiệu chỉnh:} Dựa trên phản hồi từ Tác tử Học sinh, tác tử này tiến hành sửa đổi, tinh chỉnh ngôn từ và logic để câu hỏi đạt chuẩn sư phạm cao nhất.
    \end{itemize}
\end{enumerate}

Văn bản đầu vào được xử lý bởi bốn tác tử với hai vòng lặp kiểm tra và sửa. Nhờ đó, bài toán có thể được giải quyết với chất lượng tiệm cận các LLM với chi phí thấp và có thể triển khai hoàn toàn trên hạ tầng có sẵn và bảo mật nhờ các ưu điểm của sLLM.
\section{Đóng góp của khóa luận}

Qua quá trình thiết kế, thực nghiệm và đánh giá giải pháp, khóa luận này đề xuất bốn đóng góp mang ý nghĩa cả về mặt học thuật lẫn ứng dụng thực tiễn:

\begin{itemize}
    \item \textbf{Đề xuất kiến trúc hệ thống đa tác tử chuyên biệt cho giáo dục:} Khóa luận đã thiết kế và áp dụng thành công kiến trúc đa tác tử vào các sLLM cho bài toán sinh câu hỏi trắc nghiệm. Hệ thống này có thể tự động tạo và đánh giá câu hỏi trắc nghiệm theo thang đo Bloom. Cơ chế vòng lặp kiểm tra được đề xuất để bù đắp giới hạn suy luận của sLLM.

 \item \textbf{Đánh giá thực nghiệm quy mô và toàn diện:} Khóa luận chạy các thực nghiệm trên 12 mô hình khác nhau bao gồm VLM, sLLM mã nguồn mở và các LLM thương mại cho hai bài toán sinh câu hỏi trắc nghiệm và OCR. Bên cạnh đó, dữ liệu đánh giá của tác vụ OCR được sinh tự động nhờ kỹ thuật mô phỏng nhiễu.

\item \textbf{Phát triển khung đánh giá chất lượng sư phạm:} Thay vì sử dụng các độ đo truyền thống như BLEU, ROUGE vốn dĩ không phù hợp để đánh giá các bài toán sinh dữ liệu, nghiên cứu đã thiết lập một quy trình đánh giá tự động dựa vào cơ chế LLM-as-a-Judge. Ba tiêu chí sư phạm được sử dụng là Tính giải được (Solvability), Độ bám sát phân loại nhận thức (Alignment), và Chất lượng phương án nhiễu (Distractor Quality).

 \item \textbf{Hiện thực hóa giải pháp qua ứng dụng Web Openlet:} Khóa luận đã xây dựng và triển khai ứng dụng Web nhằm ứng dụng kiến trúc đề xuất vào trong thực tế. Giáo viên có thể tải tài liệu lên để tạo bài kiểm tra và học sinh có thể tham gia bài kiểm tra đã tạo của giáo viên. Nhờ đó, nghiên cứu đã chứng minh tính khả thi của giải pháp đề xuất.

\section{Cấu trúc khóa luận}

Báo cáo sẽ trình bày một cách có hệ thống về giải pháp đề xuất, quy trình đánh gíá, kết quả và xây dựng ứng dụng Web Openlet nhằm hiện thực hóa nghiên cứu. Toàn bộ nội dung khóa được chia thành 7 chương với bố cục như sau:

 \begin{itemize}
    \item \textbf{Chương \ref{chap:mo-dau} - Mở đầu:} Đặt vấn đề, phân tích các rào cản trong các phương pháp hiện tại, từ đó đề xuất giải pháp ứng dụng sLLM kết hợp MAS, đồng thời tóm tắt các đóng góp chính của nghiên cứu.
    
    \item \textbf{Chương \ref{chap:co-so-ly-thuyet} - Cơ sở lý thuyết \& Các công trình liên quan:} Phần này cung cấp các lý thuyết nền tảng về thang đo Bloom, hệ thống đa tác tử và phương pháp LLM-as-a-Judge được sử dụng. Ngoài ra, các nghiên cứu liên quan cũng sẽ được trình bày để chỉ ra khoảng trống nghiên cứu hiện tại.
    
    \item \textbf{Chương \ref{chap:phuong-phap} - Phương pháp đề xuất:} Chương này sẽ trình bày chi tiết về kiến trúc MAS đề xuất kèm theo thiết kế của từng tác tử cũng như một số kỹ thuật mở rộng liên quan.
    
    \item \textbf{Chương \ref{chap:phuong-phap-danh-gia} - Phương pháp đánh giá:} Mục tiêu của chương này là trình bày cách khóa luận thiết lập môi trường đánh giá, mô tả bộ dữ liệu đầu vào, chiến lược tạo ảnh nhiễu và thiết lập các mô hình cơ sở dùng để đánh giá kết quả của kiến trúc đề xuất.
    
    \item \textbf{Chương \ref{chap:ket-qua} - Kết quả:} Các kết quả của thực nghiệm sẽ được đánh giá và phân tích trong chương này. Báo cáo sẽ tiến hành so sánh số liệu về năng lực sinh câu hỏi đa cấp độ, chất lượng OCR và chi phí giữa các sLLM và LLM thương mại.
    
    \item \textbf{Chương \ref{chap:thiet-ke} - Thiết kế và hiện thực ứng dụng Openlet:} Chương này mô tả quy trình xây dựng ứng dụng Web Openlet thực tế. Nội dung sẽ bao gồm sơ đồ ca sử dụng, kiến trúc của phần mềm kèm theo giới thiệu các tính năng cốt lõi của ứng dụng.
    
    \item \textbf{Chương \ref{chap:ket-luan} - Kết luận và Hướng phát triển:} Chương cuối của khóa luận sẽ tập trung vào tổng kết các kết quả quan trọng nhất mà nghiên cứu đạt được cũng như là các hạn chế đang tồn tại và một số định hướng phát triển trong tương lai.




